\subsection{Маневры и автоматический полет}

Мы уже умеем взлетать и садиться. Теперь научимся управлять движением дрона в пространстве.

\subsubsection*{1. Система координат}
Pioneer использует локальную систему координат. Представьте, что дрон находится в центре комнаты:
\begin{itemize}
    \item \textbf{X (Вперед/Назад):} Положительные значения — вперед, отрицательные — назад.
    \item \textbf{Y (Влево/Вправо):} Положительные значения — влево, отрицательные — вправо.
    \item \textbf{Z (Вверх/Вниз):} Положительные значения — вверх, отрицательные — вниз.
    \item \textbf{Yaw (Рыскание):} Поворот носа дрона (в градусах или радианах).
\end{itemize}

Команда \texttt{go\_to\_local\_point(x, y, z, yaw)} отправляет дрон в точку с координатами $(x, y, z)$ относительно места включения (старта).

\subsubsection*{2. Простой маневр: Вперед и назад}
Попробуем пролететь 1 метр вперед, подождать и вернуться обратно.

\begin{codefile}[python]{python/examples/10_01_forward_back.py}
\end{codefile}

\subsubsection*{3. Полет по квадрату}
Теперь усложним задачу. Пусть дрон облетит квадрат со стороной 1 метр.
Координаты углов:
\begin{enumerate}
    \item $(1, 0)$ — Вперед на 1 метр.
    \item $(1, 1)$ — Влево на 1 метр (теперь мы в точке $x=1, y=1$).
    \item $(0, 1)$ — Назад на 1 метр (но остаемся смещенными влево).
    \item $(0, 0)$ — Возврат в точку старта.
\end{enumerate}

\begin{codefile}[python]{python/examples/10_02_square.py}
\end{codefile}

\subsubsection*{4. Полет по кругу (Математика)}
Как заставить дрон лететь по кругу? Можно задать много точек вручную, но это долго. Лучше использовать математику!
Координаты точки на окружности вычисляются по формулам:
\[ x = R \cdot \cos(\alpha) \]
\[ y = R \cdot \sin(\alpha) \]
где $R$ — радиус, а $\alpha$ — угол.

В Python для этого используется модуль \texttt{math}.

\begin{codefile}[python]{python/examples/10_03_circle_math.py}
\end{codefile}

\begin{taskbox}[Восьмерка]
Используя формулы для лемнискаты Бернулли или просто соединив два круга, заставьте дрон пролететь по траектории <<восьмерки>>.
\end{taskbox}

\newpage
\subsection{Система контрольных упражнений}
Для проверки навыков пилотирования разработана система последовательных упражнений с автоматической проверкой точности (Grader).

\subsubsection*{Упражнение 1: Полёт вперёд}
Задача: Прямолинейный полет на заданную дистанцию (1 м) без возврата.
\begin{codefile}[python]{python/examples/exercises/ex1_forward.py}
\end{codefile}

\subsubsection*{Упражнение 2: Вперед-разворот-назад}
Задача: Полет вперед, немедленный разворот на 180 градусов и возврат в точку старта (летим <<носом>> к дому).
\begin{codefile}[python]{python/examples/exercises/ex2_forward_turn_return.py}
\end{codefile}

\subsubsection*{Упражнение 3: Строгий разворот}
Задача: Аналогично упр. 2, но с жестким контролем угла рыскания (Yaw) на всех этапах.
\begin{codefile}[python]{python/examples/exercises/ex3_forward_turn_return.py}
\end{codefile}
